% Unnumbered candidate and source audit. No original A5 closure is claimed.
% Authored by maximal_locus2, 20 November 2026, after checkpoint21.
\section*{Finite contraction certificates and a real TOP recognition test}

The finite-pair assertion below is valid over arbitrary o-minimal real
closed fields. The controlled finite-mesh argument and the ensuing
topological bundle criterion are explicitly restricted to the real
field. Neither assertion supplies a definable local bundle chart.

\textbf{1. Uniform finite certificates for nullhomotopy of compact pairs.}
Fix a definable family $(B_t,A_t)$, $t\in T$, with
\[
 \varnothing\ne A_t\subseteq B_t\subseteq R^N,
 \qquad A_t,B_t\text{ closed and bounded}.
 \tag{LC.1}
\]
Name only the coefficients of this fixed family and work in its complete
theory. There is one formula $\mathcal N(t)$, using these coefficients,
which in every model says exactly that the inclusion $A_t\hookrightarrow
B_t$ is definably nullhomotopic. Moreover, a fixed finite list of graph
formulas suffices for all its nullhomotopies; extra coefficients in a
chosen formula may depend on $t$.

Here and below a nullhomotopy means a continuous map
$H:A_t\times[0,1]\to B_t$ with $H(a,0)=a$ and $H(a,1)=z$, for one
$z\in B_t$. For a finite abstract pair $(K,L)$ with $L$ a subcomplex of
the closed complex $K$, use the standard rational realization of $K$.
For every finite graph formula $\varphi(t,u,x;c)$ form the first-order
certificate asserting that it defines a homeomorphism
\[
 \phi:|K|\longrightarrow B_t,\qquad \phi(|L|)=A_t.
 \tag{LC.2}
\]
Both continuity assertions and the marked-set equality have ordinary
finite-tuple first-order expressions. Compatible definable triangulation
gives some such certificate for every $t$, in every model of the fixed
complete theory. All those models are o-minimal: the uniform finite
component bound for each formula is a first-order sentence. Compactness,
applied to the type saying that $t\in T$ satisfies none of the
certificates, therefore gives a finite list which covers $T$ in every
model. No enumeration, decidability, or countability of the language is
being used.

For each pair in this list, the inclusion $|L|\hookrightarrow|K|$ is
definably nullhomotopic in an o-minimal real closed field if and only if
its ordinary real realization is nullhomotopic. This is the natural
homotopy-set comparison of Baro--Otero, arXiv:0810.0365v1,
Theorem~3.1 (printed p.~5), Corollary~3.4 (p.~8), with the relative
systems formulation in Corollary~3.3 (pp.~7--8) and Remark~3.5 (p.~8).
The endpoint constant may be chosen at a vertex in the appropriate
connected component of $K$: a path to that vertex appends to any
nullhomotopy. For a positive pair, choose one semialgebraic homotopy
with these endpoints over the real algebraic field. To spell out the
coefficient point, the cited comparison first gives such a finite
semialgebraic witness over a real closed field; its actual finite formula
and the homotopy identities are an ordered-field existential sentence,
which real-closed-field completeness reflects to the real algebraic
field. For a negative pair the homotopy comparison excludes a definable
witness in every model.

Label the finitely many pair types accordingly. Let $\mathcal N(t)$ be
the finite disjunction of the existence assertions (LC.2) for positive
labels. A certificate of either label has the actual pair type of
$(B_t,A_t)$, so there is no ambiguity if several certificates apply.
Transport a fixed homotopy $H_i$ by
\[
 H_{t,c}(a,s)=\phi_{t,c}
       H_i\bigl(\phi_{t,c}^{-1}(a),s\bigr).
 \tag{LC.3}
\]
This proves both the exact formula and the claimed finite witness pool.
It does not assert joint continuity when the certificate coefficients
are chosen as a function of $t$.

\textbf{2. A first-order local contraction condition.}
Let $f:X\to B\subseteq R^m$ be a fixed continuous definably proper map,
where $B$ is open, and write $F_b=f^{-1}(b)\subseteq R^N$. Its nonempty
fibres are compact. Apply Step~1 to the family of pairs
\[
 \bigl(F_b\cap\overline B_e(x),\,
       F_b\cap\overline B_d(x)\bigr),
 \qquad x\in F_b,\quad 0<d\le e.
 \tag{LC.4}
\]
Let $\mathcal N(b,x,d,e)$ denote the resulting exact predicate. The set
of points satisfying
\[
 \begin{split}
 \mathcal L(b)\ \Longleftrightarrow\
 &\forall e>0\ \exists d\in(0,e]\ \exists\eta>0\
   \forall c\in B\ \forall x\in F_c:\\[-2mm]
 &\hspace{15mm}
 \|c-b\|<\eta\quad\Longrightarrow\quad
       \mathcal N(c,x,d,e)
 \end{split}
 \tag{LC.5}
\]
is definable, in every model of the fixed named-coefficient theory.
The spatial radius $d$ is uniform over all centres $x$ in all fibres
with $\|c-b\|<\eta$; it is not merely a separate local-contractibility
assertion on each fibre. The constants $d,\eta$ may depend on $e,b$.
This proves definability of a condition on homotopies, not on
homeomorphisms.

\textbf{3. The controlled extension construction over the real field.}
We now work over $\mathbb R$. The following elementary construction will
be used twice. Suppose that a compact target $Q\subset\mathbb R^N$
has, for $j=0,\ldots,D-1$, nullhomotopies of every inclusion
\[
 Q\cap\overline B_{3r_j}(q)
       \longrightarrow Q\cap\overline B_{r_{j+1}}(q),
 \qquad q\in Q,\qquad
 0<r_0<r_1<\cdots<r_D,\qquad 3r_j\le r_{j+1}.
 \tag{LC.6}
\]
Let $K$ be a finite complex of dimension at most $D$, and give each
vertex $v$ a point $q_v\in Q$. Assume that $\|q_v-q_w\|\le r_0$
whenever $v,w$ belong to a common simplex. Let $K_0$ be a subcomplex
and let $g_0:|K_0|\to Q$ be continuous, with
\[
 g_0(\sigma)\subseteq\overline B_{r_0}(q_v)
 \quad(\sigma\in K_0,\ v\in\sigma).
 \tag{LC.7}
\]
There is an extension $g:|K|\to Q$ such that the image of each
$j$-simplex lies in $\overline B_{2r_j}(q_v)$ for every one of its
vertices. In dimension zero use $q_v$ at unprescribed vertices; the
prescribed values satisfy the same bound. Inductively, choose one vertex
$v$ of a new $j$-simplex. A boundary face has image within $2r_{j-1}$
of one of its vertices, whose assigned point is within $r_0$ of $q_v$.
Thus its whole boundary image lies in
$Q\cap\overline B_{3r_{j-1}}(q_v)$. Compose that boundary map with
the chosen nullhomotopy in (LC.6). Coning the boundary to its constant
terminal value extends it over the simplex with image inside
$\overline B_{r_j}(q_v)$. The bound for any other vertex follows by
adding $r_0\le r_j$. Prescribed simplices already satisfy the required
bound. The constructions agree on faces, so finite closed-set pasting
proves continuity. Continuity at each cone point follows from the
continuity of its nullhomotopy on a compact domain and its constant
terminal value. If the triangulations, prescribed maps and contractions
are definable, this finite construction is definable as well.

We record two consequences with constants independent of the choice of
$Q$ whenever the radii in (LC.6) work for all the target sets under
consideration. First, let $P\subset\mathbb R^N$ be compact and
triangulable, of dimension at most $N$, and suppose
$d_H(P,Q)<r_0/4$. Choose a finite triangulation of $P$ whose simplex
images have diameter less than $r_0/4$, and at each vertex choose
$q_v\in Q$ within $r_0/4$. The vertex condition above holds. With
$D=N$, the extension construction, before changing reference vertices,
places each simplex image within $r_N$ of one of its assigned points.
Consequently it gives a map $g:P\to Q$ with
\[
             \sup_{x\in P}\|g(x)-x\|<2r_N.
 \tag{LC.8}
\]

Second, let $P$ be a compact triangulable space of dimension at most $N$
and let $u,v:P\to Q$ satisfy $\|u-v\|_\infty<r_0/4$. Subdivide a
finite triangulation of $P$ until $u$ maps every simplex to a set of
diameter less than $r_0/4$. Triangulate its product with $[0,1]$ using
only the vertices at the two ends. Assign $u(a)$ to a bottom vertex
and $v(a)$ to its top counterpart. Points assigned to vertices in one
product simplex differ by less than $3r_0/4$. Prescribe $u,v$ on the
bottom and top subcomplexes. On either prescribed simplex the image is
within $r_0$ of every assigned vertex value, by the triangle inequality.
Apply the relative construction with $D=N+1$. Every homotopy value over
$x\in P$ differs from $u(x)$ by less than $3r_{N+1}$: its controlling
vertex contributes at most $2r_{N+1}$, and its assigned point differs
from $u(x)$ by at most $r_0/2$. Thus, if $r_{N+1}<e/6$, this is a
homotopy from $u$ to $v$ each of whose tracks has diameter less than $e$.

The finite meshes just used follow from ordinary compactness and uniform
continuity on the real field. A standard finite subdivision need not have
infinitesimal mesh over a non-Archimedean field. Accordingly this argument
neither proves an all-field assertion nor bounds its number of simplices
as the desired accuracy tends to zero.

\textbf{4. From (LC.5) to strong regularity over $\mathbb R$.}
Suppose all fibres are nonempty, $f$ is open, and (LC.5) holds at $b$.
Properness and openness imply Hausdorff continuity of the fibres at $b$.
For completeness, restrict to a compact target neighbourhood of $b$;
its inverse image is compact. Upper semicontinuity follows by taking a
convergent subsequence of any putative sequence of offending fibre
points. Lower semicontinuity follows from openness at each point of
$F_b$ and a finite cover of that compact fibre by small balls. These two
statements are exactly Hausdorff continuity.

Given a desired homotopy size $e>0$, repeatedly apply (LC.5), working
backwards from $r_{N+1}<e/6$. At each stage choose
$3r_j$ smaller than the radius supplied for the target radius
$r_{j+1}$, and also $r_j<r_{j+1}/4$. Intersect the finitely many target
neighbourhoods so that (LC.6) holds for every nearby fibre, including
$F_b$. Thus the second consequence of Step~3 supplies a uniform
threshold $a>0$ such that two maps into any of these fibres which are
$a$-close are joined by an $e$-small homotopy.

Choose $\mu<\min\{e,a/3\}$. Use a second finite backward chain ending
with $r_N<\mu/2$, and again shrink the target neighbourhood. Hausdorff
continuity and the first consequence of Step~3 give, for every nearby
$c$, maps
\[
 f_c:F_c\to F_b,\qquad g_c:F_b\to F_c,
 \qquad \|f_c-\mathrm{id}\|_\infty<\mu,\quad
        \|g_c-\mathrm{id}\|_\infty<\mu.
 \tag{LC.9}
\]
These norms compare ambient points; the displayed identities are not
asserted to have the same codomains. The compositions differ from their
respective identities by less than $2\mu<a$. The first chain therefore
gives $e$-small homotopies of these compositions to the identities, taking
place in the respective fibres. This is strong regularity at $b$.
No continuity of the choices $f_c,g_c$ in $c$ has been asserted.

\textbf{5. An exact real TOP criterion for closed manifold fibres of
dimension at least five.}
Fix $n\ge5$. Let $M_n$ be the set of $b$ with nonempty fibre a closed
topological $n$-manifold. It is definable: a finite Hardt partition gives
finitely many compact fibre homeomorphism types, and one retains the
manifold types. Let $O$ be the definable set of $b$ such that $f$ is open
at every point of $F_b$; an epsilon--delta formula expresses this
condition. Then the exact real topological bundle locus with closed
$n$-manifold fibre is
\[
               U_{\mathrm{TOP},n}=\operatorname{Int}_B
                         \bigl(M_n\cap O\cap\mathcal L\bigr).
 \tag{LC.10}
\]

To prove necessity, a bundle chart gives openness and the stated fibre
type on its open base neighbourhood. Restrict it to a smaller compact
base neighbourhood. Its evaluation and inverse are uniformly continuous.
A compact triangulable fibre is uniformly locally contractible: finitely
many small open stars in a triangulation provide contractions inside any
prescribed neighbourhood scale, and a Lebesgue number gives a uniform
input radius. Transporting these contractions through the compact chart
proves (LC.5) throughout a neighbourhood of the given point. If the chart
is merely topological, the resulting inclusion nullhomotopies still
give the same truth values in (LC.4), by the finite-pair comparison in
Step~1. Thus a topological bundle point lies in the right-hand interior.

For sufficiency, work on any open set contained in that right-hand side.
Step~4 gives strong regularity everywhere. We spell out its conversion
to the hypothesis which actually recognizes bundles. For fixed $b$ and
an arbitrarily small open cover $\alpha$ of the compact manifold $F_b$,
the Chapman--Ferry alpha-approximation theorem gives a cover $\beta$
of $F_b$. Choose a Lebesgue number for $\beta$. Strong regularity gives
nearby-fibre maps and homotopies sufficiently small that $f_c$ is a
$\beta$-equivalence: an ambient-small source homotopy remains small
after $f_c$, because $f_c$ itself moves ambient points only slightly.
There is no boundary condition to arrange for closed fibres. The
alpha-approximation theorem replaces $f_c$ by an $\alpha$-close
homeomorphism. It follows, by choosing both scales small, that nearby
fibres admit arbitrarily small ambient-displacement homeomorphisms.
This is complete regularity.

Chapman--Ferry Proposition~7.1 then gives local triviality, since the
base is locally finite-dimensional and these compact triangulable fibre
homeomorphism groups are locally contractible. For the latter assertion
one can use Gauld's original \emph{Local contractibility of spaces of
homeomorphisms}, Compositio Mathematica 32 (1976), 3--11,
Theorem~(1) and its Corollary on printed p.~3: it applies to every finite
polyhedron. Its limiting isotopy is not asserted here to be definable.
The actual Chapman--Ferry source is
\emph{Approximating homotopy equivalences by homeomorphisms},
American Journal of Mathematics 101 (1979), 583--607:
alpha-approximation is stated on printed p.~583; the Bundle Theorem is
on p.~584; complete regularity, Proposition~7.1, strong regularity and
the start of the bundle proof are on p.~605, with its completion on
p.~606. The threshold is $n\ge5$, including dimension five. The source
is topological, not semialgebraic. The argument here uses its
strong-to-complete-regularity step directly and does not confuse strong
regularity with complete regularity.

\textbf{6. Exact remaining limitations.}
Formula (LC.10) concerns the real TOP locus and closed manifold fibres
of dimension at least five. It does not prove that a topological chart
of a definable map can be replaced by a definable chart. It does not
settle fibres of dimensions three or four, arbitrary compact polyhedral
fibres, or the all-real-closed-field A5 assertion. In particular the
finite pool (LC.3) is a pool of nullhomotopies of ball pairs, not a pool
of the exact angular homeomorphisms in A5.019. No implication between
those two pools is assumed.

The source trace and original locators above were checked directly on
20 November 2026. Baro--Otero's complete v1 original was opened at
\url{https://arxiv.org/pdf/0810.0365}; Chapman--Ferry's complete original
was opened at
\url{https://webhomes.maths.ed.ac.uk/~v1ranick/papers/chapferr.pdf}.
The latter's printed pp.~583--584 and 605--606 were also read as page
images. The downloaded complete external original is kept outside the
authored project, under \path{/tmp/A5_originals_chapman_ferry/}.
Gauld's complete original was opened at
\url{https://www.numdam.org/item/CM_1976__32_1_3_0.pdf}.
